\documentclass[10pt,twocolumn]{article}
\usepackage[a4paper,margin=1.8cm]{geometry}
\usepackage{amsmath,amssymb,graphicx,booktabs,multirow,subcaption,algorithm2e,xcolor}
\title{Stress Test Two: A Two-Column Paper}
\author{C. Columns}
\date{}
\begin{document}
\maketitle
\begin{abstract}
Two-column layouts are the norm in engineering venues. This paper mixes running text, equations that
fill a narrow column, a figure spanning both columns, a wide table and an algorithm.
\end{abstract}

\section{Introduction}
Neural networks compute $y = \sigma(Wx + b)$ where $\sigma$ is a nonlinearity such as the rectified
linear unit $\mathrm{ReLU}(z) = \max(0, z)$. Training minimises the empirical risk
\begin{equation}
  \mathcal{L}(\theta) = \frac{1}{N}\sum_{i=1}^{N} \ell\big(f_\theta(x_i), y_i\big) + \lambda \|\theta\|_2^2 ,
  \label{eq:risk}
\end{equation}
with gradient descent $\theta_{t+1} = \theta_t - \eta \nabla_\theta \mathcal{L}(\theta_t)$.
The learning rate $\eta$ matters a great deal.\footnote{See the appendix for the schedule we used.}
Results are summarised in Table~\ref{tab:main} and Figure~\ref{fig:wide}.

The softmax of a vector $z \in \mathbb{R}^K$ is
\begin{equation}
  \operatorname{softmax}(z)_k = \frac{e^{z_k}}{\sum_{j=1}^{K} e^{z_j}} .
\end{equation}
Attention combines queries, keys and values:
\begin{equation}
  \mathrm{Attn}(Q,K,V) = \operatorname{softmax}\!\left(\frac{QK^\top}{\sqrt{d_k}}\right) V .
\end{equation}

\section{Method}
We use the procedure in Algorithm~\ref{alg:train}. Each step samples a batch $B \subset \{1,\dots,N\}$ of size $|B| = 64$.

\begin{algorithm}[t]
\caption{Training loop}\label{alg:train}
\KwIn{data $\{(x_i,y_i)\}_{i=1}^N$, rate $\eta$}
\For{$t = 1$ \KwTo $T$}{
  sample batch $B$\;
  $g \leftarrow \frac{1}{|B|}\sum_{i\in B}\nabla \ell_i$\;
  $\theta \leftarrow \theta - \eta g$\;
}
\KwOut{$\theta$}
\end{algorithm}

The model has three parts. First, an encoder maps the input to a hidden state. Second, a decoder reads the
hidden state. Third, a small head predicts the label. Each part is trained jointly, and we report the
average over five random seeds. This paragraph is deliberately long so that it flows from one column into
the next one, which is where many converters lose the reading order or mix the two columns together line
by line. A careful reader should see this sentence exactly once and in the right place.

\begin{figure*}[t]
  \centering
  \begin{subfigure}{0.32\textwidth}\includegraphics[width=\linewidth]{lines.pdf}\caption{Lines}\end{subfigure}\hfill
  \begin{subfigure}{0.32\textwidth}\includegraphics[width=\linewidth]{bars.pdf}\caption{Bars}\end{subfigure}\hfill
  \begin{subfigure}{0.32\textwidth}\includegraphics[width=\linewidth]{scatter.pdf}\caption{Scatter}\end{subfigure}
  \caption{Three panels side by side spanning both columns. Panel (a) shows periodic signals, (b) a bar
  chart and (c) a scatter plot.}
  \label{fig:wide}
\end{figure*}

\begin{table*}[t]
  \centering
  \caption{Main results on six benchmarks. Best numbers are in bold.}
  \label{tab:main}
  \begin{tabular}{lcccccccc}
    \toprule
    \multirow{2}{*}{Method} & \multicolumn{2}{c}{Vision} & \multicolumn{2}{c}{Language} & \multicolumn{2}{c}{Audio} & \multirow{2}{*}{Avg.} & \multirow{2}{*}{Params} \\
    \cmidrule(lr){2-3}\cmidrule(lr){4-5}\cmidrule(lr){6-7}
     & Top-1 & Top-5 & BLEU & PPL & WER & CER & & \\
    \midrule
    Baseline & 71.2 & 90.1 & 27.3 & 18.4 & 9.1 & 3.2 & 36.6 & 25M \\
    Wide & 73.5 & 91.0 & 28.1 & 17.2 & 8.7 & 3.0 & 37.6 & 60M \\
    Deep & 74.0 & 91.8 & \textbf{29.4} & 16.9 & 8.5 & 2.9 & 38.3 & 80M \\
    Ours & \textbf{76.8} & \textbf{93.2} & 29.1 & \textbf{15.8} & \textbf{7.9} & \textbf{2.6} & \textbf{39.2} & 45M \\
    \bottomrule
  \end{tabular}
\end{table*}

\section{Results}
Our model improves accuracy by $2.8$ points while using fewer parameters than \emph{Deep}. The perplexity
drops from $16.9$ to $15.8$, a relative gain of $\approx 6.5\%$. Error rates fall as well.

\begin{table}[h]
  \centering
  \caption{Ablation with grid lines.}
  \begin{tabular}{|l|r|r|}
    \hline
    Variant & Acc. & $\Delta$ \\ \hline
    Full model & 76.8 & -- \\ \hline
    no attention & 74.1 & $-2.7$ \\ \hline
    no dropout & 75.9 & $-0.9$ \\ \hline
  \end{tabular}
\end{table}

\section{Conclusion}
We presented a stress test. The inequality $\mathbb{E}[\|g\|^2] \le G^2$ bounds the gradient norm, and the
bound $\mathcal{O}(1/\sqrt{T})$ follows. Future work includes $k$-fold validation with $k \in \{5, 10\}$.
\end{document}
